\documentclass[10pt]{amsart}
\usepackage{amssymb,amsmath,amsthm,amsfonts,amsopn,url}
\usepackage{amscd,amssymb,amsopn,amsmath,amsthm,graphics,amsfonts,enumerate,verbatim,calc}
\usepackage[all]{xy}
\theoremstyle{plain}
\newtheorem{thm}{Theorem}[section]
\newtheorem*{mt*}{Theorem (Huneke-Lyubeznik)}
\newtheorem*{cj*}{Conjecture}
\newtheorem*{nt*}{Notations}
\newtheorem*{stfac}{Base change formulas for depth}
\newtheorem{prop}[thm]{Proposition}
\newtheorem{lemma}[thm]{Lemma}
\newtheorem{cor}{Corollary}
\newtheorem{rem}{Remark}
\newtheorem{example}{Example}[section]
\theoremstyle{definition}
\newtheorem*{definition}{Definition (The graph $\Theta_R$)}
\newcommand{\ideal}[1]{\mathfrak{#1}}
\newcommand{\M}{\ideal{M}}
\newcommand{\m}{\ideal{m}}
\newcommand{\n}{\ideal{n}}
%\newcommand{\o}{\ideal{o}}
\newcommand{\p}{\ideal{p}}
\newcommand{\q}{\ideal{q}}
\newcommand{\I}{\ideal{I}}
\newcommand{\J}{\ideal{J}}
\newcommand{\ffunc}[1]{\mathrm{#1}}
\newcommand{\func}[1]{\mathrm{#1} \,}
\newcommand{\depth}{\func{depth}}
\newcommand{\depthsb}{\ffunc{depth}}
\newcommand{\height}{\func{ht}}
\newcommand{\grade}{\func{grade}}
\newcommand{\ara}{\func{ara}}
\newcommand{\gldim}{\func{gl} \dim}
\newcommand{\Tor}{\func{Tor}}
\newcommand{\Hom}{\func{Hom}}
\newcommand{\Ass}{\func{Ass}}
\newcommand{\Ann}{\func{Ann}}
\newcommand{\im}{\func{im}}
\newcommand{\V}{\func{V}}
\newcommand{\rad}{\func{rad}}
\newcommand{\Supp}{\func{Supp}}
\newcommand{\Spec}{\func{Spec}}
\newcommand{\Ker}{\func{Ker}}
\newcommand{\Coker}{\func{Coker}}
\newcommand{\pdp}[1]{\ffunc{ph. depth}_{#1}}
\newcommand{\arrow}[1]{\stackrel{#1}{\rightarrow}}
\newcommand{\PP}{{\mathbb P}}
\newcommand{\QQ}{{\mathbb Q}}
\newcommand{\NN}{{\mathbb N}}
\newcommand{\ZZ}{{\mathbb Z}}
\newcommand{\CC}{{\mathbb C}}


\title[]{local cohomology module with infinitely many associated primes over ramified regular local rings of the least possible dimensions}

\author[]{Rajsekhar Bhattacharyya}

\address{Dinabandhu Andrews College, Garia, Kolkata 700084, India}

\email{rbhattacharyya@gmail.com}

\thanks{}

\keywords{}

\subjclass[2010]{13D45}
\begin{document}

\begin{abstract}
In \cite{Ma} we have observed the example of local cohomology module with infinitely many associated primes over a 12 dimensional ramified regular local ring. In this paper we construct examples of local cohomology modules with infinitely many associated primes over ramified regular local rings of two different ring dimensions. The result of this paper is very much dependent on the result (\cite{SS}, Theorem 2.6). 
\end{abstract}

\maketitle




\section{}
In this paper we construct examples of a local cohomology module with infinitely many associated primes over ramified regular local rings of the least possible dimensions. Here we exhibit two examples for two different ring dimensions. Here we follow similar construction of that of \cite{Ma}.\\

{\bf 1st Example}\\

For the 1st example we proceed as follows: Let $Z_{(2)}$ be the DVR which is obtained from the ring of integers after inverting every integers other than 2 and $F_2$ is the prime field of characteristic 2. Let $I=(x_1) \subset Z_{(2)}[[x_1]]$ be an ideal.
It follows that
$0\neq H^1_I (Z_{(2)}[[x_1]]) =Z_2[[x_1]]_{(x_1)}/Z_2[[x_1]]$ and $H^2_I (Z_2[[x]]) = 0.$

Next, we recall the example of a local cohomology module of a hypersurface with infinitely many
associated primes from (\cite{SS}, Remark 4.2) (followed by a
completion map):

$H^2_{(y1,y2)} (F_2[[y_1,\ldots, y_5]]/(y^2_1 y^2
_3  + y_1y_2y_3y_4y_5 + y^2_2 y^2_4 ))
$
has infinitely many associated primes. It is to be noted the above example follows at once from (\cite{SS}, Theorem 2.6). 

Now we set $S = Z_{(2)}[[x, y]] := Z_2[[x_1, y_1,\ldots, y_5]]$ and let
$R = S/(2 + y^2_1 y^2
_3 + y_1y_2y_3y_4y_5 + y^2_2 y^2_4)$.
Note that R is a complete ramified regular local ring. Finally, we let
$J = IS + (y1, y2)S\subset S$.

We claim that $H^3_J (R)$ has infinitely many associated primes. We can prove as follows:

It is well-known that for any ring $A$ and any ideal $\alpha\subset A$, we
have $H^{i+1}
_{\alpha+(z)}(A[[z]]) = H^i_\alpha (A) \otimes_Z H^1_{(z)}(Z[[z]])$. This gives:
$H^3_J (S) = H^1_I (Z_{(2)}[[x_1]])\otimes_Z H^2_{(y_1,y_2)}(Z[[y]])$ %= E_{F_2[[x]]}(F_2)\otimes_{F_2} H^2_{(y1,y2)}(F_2[[y]])$,
and
$H^4_J (S) = H^2_I (Z_{(2)}[[x]])\otimes_Z H^2_{(y1,y2)}(Z[[y]]) = 0$

\newpage




Denote $f (y) := y^2_1 y^2
_3  + y_1y_2y_3y_4y_5 + y^2
_2 y^2
_4 $ and consider part of the long exact sequence

$H^3_J (S)\stackrel{(2+f (y))}{\rightarrow} H^3_J (S)\rightarrow H^3_J (R)\rightarrow H^4_J (S) = 0$.

It follows that $H^3_J (R)$ is the cokernel of the multiplication by $(2 + f (y))$ map on $H^3_J (S)$. But since
$H^3_J (S) = H^1_I (Z_{(2)}[[x_1]])\otimes_Z H^2_{(y_1,y_2)}(Z[[y]])$. Now image of multiplication map is $(2 + f (y))(H^1_I (Z_{(2)}[[x_1]])\otimes_Z H^2_{(y_1,y_2)}(Z[[y]]))= H^1_I (Z_{(2)}[[x_1]])\otimes_Z (2 + f (y))(H^2_{(y_1,y_2)}(Z[[y]]))$. 
So $H^3_J (R)= \Coker (2 + f (y))= H^1_I (Z_{(2)}[[x_1]])\otimes_Z \Coker (H^2_{(y_1,y_2)}(Z[[y]]\stackrel{(2+f (y))}{\rightarrow}H^2_{(y_1,y_2)}(Z[[y]])= H^1_I (Z_{(2)}[[x_1]])\otimes_Z (H^2_{(y_1,y_2)}(\Coker (Z[[y]]\stackrel{(2+f (y))}{\rightarrow}Z[[y]])=H^1_I (Z_{(2)}[[x_1]])\otimes_Z (H^2_{(y_1,y_2)}(F_2[[y]]/f (y)) $, where we used that $\Coker H^2_{(y1,y2)}(Z[[y]]) \stackrel{(2+f (y))}{\rightarrow} H^2_{(y1,y2)}(Z[[y]]) = H^2_{(y1,y2)}(F_2[[y]]/f (y))$ since the
cohomological dimension of $(y_1, y_2)$ is two. Now $H^2_{(y1,y2)}(F_2[[y]]/f (y))$ is an $F_2$-vector space and hence we can write:

$H^3_J (R)= (H^1_I (Z_2[[x_1]])\otimes_Z F_2)\otimes_{F_2} H^2_{(y1,y2)}(F_2[[y]])$ and in particular $H^3_J (S)$ is annihilated by $(2)$.

Since $H^3_J (R)$ is annihilated by $(2)$, it is enough to show that $H^3_J (R)$ has infinitely many associated primes as a module over $R/(2)$. But clearly we have $R/(2)= F_2[[y_1,\ldots, y_5]]/(y^2_1 y^2_3  + y_1y_2y_3y_4 + y^2_2 y^2_4) [[x_1]]$.

By the explicit description of $H^3_J (R)$ above, for one fixed asssociated prime $P$ of $H^1_I (Z_2[[x_1]])\otimes_Z F_2$ over $F_2[[x_1]]$ it is clear that for each of the infinitely many associated
prime $Q$ of the local cohomology module in (*), $Q + P$ is an associated prime of $H^3_J (R)$
over $R/(2)$. In particular, $H^3_J (R)$ has infinitely many associated primes. 

So in this example dimension of $R$ is 6.\\

{\bf 2nd Example}\\
 
From (\cite{SS}, Theorem 2.6), we can get another example of local cohomology module (followed by a completion map) $H^2_{(y_1,y_2)} (F_2[[y_1,\ldots, y_4]]/(y^2_1 y^2_3  + y_1y_2y_3y_4 + y^2_2 y^2_4 ))$ with infinitely many associated primes. In fact this follows directly from (\cite{SS}, Theorem 2.6), after we set $F_2, y_1, y_2, y_3, y_4, 2, 1, 1, 1$  for $A_0, u, v, x, y, d, r_0, r_1, r_2$ of Theorem 2.6 respectively. Now using this example, we can proceed exactly in the same way above, and we can get a counter example of local cohomology with infinitely many associated primes in a ramified regular local ring of dimension 5, but this contradicts Boix's recent work. Since we can not get examples of infinite asociated primes in dimension less that 4 due to \cite{Mar}, we can not go down bellow dimension 5. 

So in this example dimension of $R$ is 5.



%%%%%%%%%%%%%%%%%%%%%%%%%%%%%%%%%%%%%%%%%%%%%%%%%%%%%%%%%%%%%%%%%%%%%%%%%%%%%%%%%%%%%%%%%%%%%%%%%%%%%%%%%%%%%%%%%%%%%%%%%%%%%%%%%%%%%%%%
\begin{thebibliography}{99}
\bibitem{Boix} Alberto F.Boix, {\it Associated primes of local cohomology modules for five dimensional regular local rings}, DOI: 10.13140/RG.2.2.11073.52322

\bibitem{Ma}
%\bibitem[Mar]{Mar}
Linquan. Ma, \emph{Infinitely many associated primes of local cohomology modules of ramified regular local rings}, arXiv:  arXiv:2604.10964v2

\bibitem{Mar}
T. Marley, \emph{The associated primes of local cohomology modules over rings of small dimension}, Manuscripta Mathematica, Volume 104, pages 519-525 (2001).

\bibitem{SS}
%\bibitem[SS]{SS}
Anurag K. Singh and Irena Swanson, {\it Associated primes of local cohomology modules and of Frobenius powers}, Int. Math. Res. Not., no. 33, 1703-1733, 2004.
\end{thebibliography}


\end{document}
